\documentclass[12pt]{article}
\usepackage[utf8]{inputenc}
\usepackage{graphicx}
\usepackage{hyperref}
\usepackage{amssymb}
\usepackage{amsmath}
\usepackage[T1]{fontenc}
\usepackage{lmodern}
\usepackage{textcomp}
\usepackage{pgfplots}
\usepackage{pgfplotstable}
\usepackage{longtable}
\usepackage[capitalise]{cleveref}
\usepackage{caption}
\pgfplotsset{compat=1.18}
\captionsetup[table]{position=top,justification=raggedright,singlelinecheck=false}
\captionsetup[longtable]{position=top,justification=raggedright,singlelinecheck=false}
\captionsetup[figure]{justification=raggedright,singlelinecheck=false}
\setlength\LTleft{0pt}
\setlength\LTright{0pt}
\setlength\LTcapwidth{\textwidth}
\pgfplotstableset{
  show variable once/.style={
    columns/variable/.style={
      string type,
      column name=Variable,
      postproc cell content/.code={
        \pgfkeysgetvalue{/pgfplots/table/@cell content}{\cellcontent}
        \edef\currentvar{\cellcontent}
        \ifx\currentvar\lastvar
          \pgfkeyssetvalue{/pgfplots/table/@cell content}{}
        \else
          \ifnum\pgfplotstablerow>0
            \pgfkeyssetevalue{/pgfplots/table/@cell content}{\noexpand\hline\space\cellcontent}
          \fi
          \xdef\lastvar{\currentvar}
        \fi
      }
    }
  }
}
\def\lastvar{}
\setlength{\parskip}{0.8em}
\setlength{\parindent}{0pt}
%\renewcommand{\thesubsection}{\thesection.\alph{subsection}}

\title{\Large\bfseries ASSIGNMENT 4: GENERAL LINEAR MODEL — MODEL SELECTION}
\author{Chenghao Wang}
\date{\today}

\begin{document}

\maketitle

Mean model estimation and covariance model estimation are interwoven.
Therefore, mean model selection depends on a correctly-specified
covariance model; covariance model selection also depends on a
correctly-specified mean model.

One strategy is that we first select the mean model and then select
the covariance model. However, 
as is described in Assignment 3,
the dataset contains 16 time points,
which means $16 \times 15 / 2 = 120$ parameters
in the unstructured correlation matrix, 
resulting in a high computational cost.
Thus, in this assignment, we adopt a different strategy.
We first select the covariance model;
then we select the mean model; lastly, we verify
the preliminary covariance structure is suitable.

R package \texttt{nlme} and \texttt{mmrm} were used.

\section{Covariance Model Selection}\label{sec:Covariance-Model-Selection}

\subsection{Candidate covariance structures}\label{subsec:candidate-covariance-structures}

To find the covariance structure that provides the most adequate fit,
we fit saturated mean models (one parameter per group-bytime cell) adjusting
for three covariates (age at randomization, gender, and ethnic group),
paired with five correlation structures respectively,
which are Compound Symmetry, AR(1), Exponential, Toeplitz, and Ante-dependence (ANTE).
For a description of ANTE, see
\href{https://cran.r-project.org/web/packages/mmrm/vignettes/covariance.html#homogeneous-ad-and-heterogeneous-ante-dependence-adh}{here}.
In addition, we assumed homogeneous variances at this step,
while the variance-covariance structures were evaluated
after the selction of correlation structure.
The models were estimated using REML.
The correlation structures comparison 
is shown in \cref{tab:corr-model-compare}.

\begin{table}[htbp]
\caption{Model comparison for candidate correlation structures}
\label{tab:corr-model-compare}
\centering
\pgfplotstabletypeset[
  col sep=comma,
  columns={model,dof,logLik,AIC},
  columns/model/.style={column name=Model,string type},
  columns/dof/.style={column name=DF,column type=r},
  columns/logLik/.style={
    column name=logLik,
    column type=r,
    fixed,
    precision=3
  },
  columns/AIC/.style={
    column name=AIC,
    column type=r,
    fixed,
    precision=3
  },
  every head row/.style={before row=\hline,after row=\hline},
  every last row/.style={after row=\hline}
]{../hw4_code/output/model_compare.csv}
\end{table}



\subsection{Model comparison}

ANTE shows the lowest AIC, indicating that the model
with ANTE correlation structure exhibits the most
adequate fit among the five candidates.
Although we selected Exponential correlation structure
in Assignment 3 due to unequally spaced time points,
yet we also reported that the correlation with fixed lag
slightly decreases over time, which can
be accommodated by ANTE.

Next, assuming the mean structure identical to the one
in \cref{subsec:candidate-covariance-structures} and
ANTE correlation structure, we fit the models
with four variance-covariance structures respectively,
which are listed in \cref{tab:var-model-compare}.
Similary, the models were estimated using REML.
Incidentally, ANTE is available in \texttt{mmrm}
but not in \texttt{nlme}. And the diverse variance structures
that we inspected in Assignment 3
are only supported by \texttt{nlme} but not \texttt{mmrm}.
Besides, group-specific covariance is supported by
\texttt{mmrm}, while \texttt{nlme} only supports
group-specific variance and does not support 
group-specific correlation.

\begin{table}[htbp]
\caption{Model comparison for candidate variance-covariance structures under ANTE}
\label{tab:var-model-compare}
\centering
\pgfplotstabletypeset[
  col sep=comma,
  columns={model,dof,logLik,AIC},
  columns/model/.style={column name=Model,verb string type},
  columns/dof/.style={column name=DF,column type=r},
  columns/logLik/.style={
    column name=logLik,
    column type=r,
    fixed,
    precision=3
  },
  columns/AIC/.style={
    column name=AIC,
    column type=r,
    fixed,
    precision=3
  },
  every head row/.style={before row=\hline,after row=\hline},
  every last row/.style={after row=\hline}
]{../hw4_code/output/aic_compare_mmrm_var.csv}
\end{table}

As shown in \cref{tab:var-model-compare},
Hetergeneous variance with group-specific ANTE covariance
provides the most adequate fit.

\section{Mean Model Selection}\label{sec:Mean-Model-Selection}

\subsection{Candidate mean structures}

\begin{table}[htbp]
\caption{Model comparison for candidate mean structures}
\label{tab:mean-model-compare}
\centering
\pgfplotstabletypeset[
  col sep=comma,
  columns={model,dof,logLik,AIC,BIC},
  columns/model/.style={column name=Model,string type},
  columns/dof/.style={column name=DF,column type=r},
  columns/logLik/.style={
    column name=logLik,
    column type=r,
    fixed,
    precision=3
  },
  columns/AIC/.style={
    column name=AIC,
    column type=r,
    fixed,
    precision=3
  },
  columns/BIC/.style={
    column name=BIC,
    column type=r,
    fixed,
    precision=3
  },
  every head row/.style={before row=\hline,after row=\hline},
  every last row/.style={after row=\hline}
]{../hw4_code/output/mean_aic_compare.csv}
\end{table}

Using group-specific heterogeneous ANTE covariance structure, 
We fit models with four candidate mean structures,
which are linear trend, linear spline (knot = 60, selected
based on the mean response trajectories plot in Assignment 3),
quadratic trend, and saturated mean model. 
The baseline mean responses of each treatment group
adjusting for baseline age, gender, and ethnic group
were constrained to be equal under
the linear and the linear spline models,
by excluding the main effect
of each treatment group.
Under the quadratic model, such constraint was not
imposed, while time was centered to reduce collinearity.
The models were
estimated using ML. The mean structure comparison
is illustrated in \cref{tab:mean-model-compare}.

\subsection{Model comparison}

As shown in \cref{tab:mean-model-compare},
The saturated mean model exhibits the lowest AIC by a very
large margin,
suggesting that the gain in fit from this model
is worth the extra parameters.
However, the saturated mean model gives the highest BIC,
indicating that this model is not preferred,
if we impose a stronger complexity penalty.
Additionally, the linear spline model (knot = 60)
shows the second lowest AIC and the lowest BIC,
indicating an acceptable fit.
Thus we selected the linear spline model (knot = 60)
as the final model, because it acheives a more
parsimonious representation and is more interpretable.
Nonetheless, a further verifaction on the covariance structure
is desirable, which is shown in \cref{subsec:Preliminary-covariance-structure-verification}.

\subsection{Preliminary covariance structure verification}\label{subsec:Preliminary-covariance-structure-verification}

Finally, we re-evaluated the covariance structures
using the linear spline model (knot = 60). The models were
estimated using REML. As shown in \cref{tab:lsp-corr-model-compare},
ANTE provides the most adequate fit
among the linear spline models with homogeneous variance
and various correlation structures.
Furthermore, under ANTE,
the model with heterogeneous variance with group-specific covariance
provides the most adequate fit (\cref{tab:lsp-var-model-compare}).
These results confirm that the preliminary choice of
covariance structure remains the most acceptable one
under the linear spline model.

\begin{table}[htbp]
\caption{Correlation structure comparison under the linear spline model}
\label{tab:lsp-corr-model-compare}
\centering
\pgfplotstabletypeset[
  col sep=comma,
  columns={model,dof,logLik,AIC},
  columns/model/.style={column name=Model,string type},
  columns/dof/.style={column name=DF,column type=r},
  columns/logLik/.style={
    column name=logLik,
    column type=r,
    fixed,
    precision=3
  },
  columns/AIC/.style={
    column name=AIC,
    column type=r,
    fixed,
    precision=3
  },
  every head row/.style={before row=\hline,after row=\hline},
  every last row/.style={after row=\hline}
]{../hw4_code/output/lsp_cov_aic_compare.csv}
\end{table}

\begin{table}[htbp]
\caption{Variance structure comparison under the linear spline model}
\label{tab:lsp-var-model-compare}
\centering
\pgfplotstabletypeset[
  col sep=comma,
  columns={model,dof,logLik,AIC},
  columns/model/.style={column name=Model,string type},
  columns/dof/.style={column name=DF,column type=r},
  columns/logLik/.style={
    column name=logLik,
    column type=r,
    fixed,
    precision=3
  },
  columns/AIC/.style={
    column name=AIC,
    column type=r,
    fixed,
    precision=3
  },
  every head row/.style={before row=\hline,after row=\hline},
  every last row/.style={after row=\hline}
]{../hw4_code/output/lsp_var_aic_compare.csv}
\end{table}



\section{Final Model}

\subsection{Model specification}

The final model formulation is as follows:

\begin{multline}
\label{eq:final-model}
Y_{i} = \beta_{0} + I(TG_i = bud)\,visit_{i}\beta_{visit \times bud}
+ I(TG_i = bud)\,(visit_{i} - 60)_{+}\beta_{visit \times bud}^{sp}\\
+ I(TG_i = ned)\,visit_{i}\beta_{visit \times ned} 
+ I(TG_i = ned)\,(visit_{i} - 60)_{+}\beta_{visit \times ned}^{sp} \\
{}+ I(TG_i = plbo)\,visit_{i}\beta_{visit \times plbo}
+ I(TG_i = plbo)\,(visit_{i} - 60)_{+}\beta_{visit \times plbo}^{sp}\\
+ age\_rz_{i}\,\beta_{age\_rz}
+ I(gender_{i} = m)\,\beta_{gender}\\
{}+ I(ethnic_{i} = b)\,\beta_{ethnic:black}
+ I(ethnic_{i} = h)\,\beta_{ethnic:hispanic}\\
{}+ I(ethnic_{i} = o)\,\beta_{ethnic:other}
+ \varepsilon_{i}\\
\varepsilon_{i} \sim N(0,\Sigma_{TG_i}),
\end{multline}

where $(x)_{+} = \max(x, 0)$, and
$\Sigma_{TG_i}$ is a group-specific hetergeneous ANTE covariance matrix.
The reason for choosing this model is discussed in 
\cref{sec:Covariance-Model-Selection} and 
\cref{sec:Mean-Model-Selection}.


\subsection{Parameter estimates}

The fixed effect estimates, model-based standard errors,
t-statistics, degrees of freedom (estimated via Between-Within method),
p-value, and 95\% confidence intervals are reported
in \cref{tab:parameter-estimates}.

\begin{table}[htbp]
\caption{Parameter estimates of the final model}
\label{tab:parameter-estimates}
\centering
\resizebox{\textwidth}{!}{%
\pgfplotstabletypeset[
  col sep=comma,
  header=false,
  skip first n=1,
  text indicator=",
  columns={0,1,2,3,4,5,6},
  columns/0/.style={column name=Effect,verb string type},
  columns/1/.style={column name=Estimate,column type=r,fixed,precision=3},
  columns/2/.style={column name=SE,column type=r,fixed,precision=3},
  columns/3/.style={column name=t,column type=r,fixed,precision=3},
  columns/4/.style={column name=DF,column type=r,fixed,precision=1},
  columns/5/.style={column name=p-value,column type=r,string type},
  columns/6/.style={column name=95\% CI,string type},
  every head row/.style={before row=\hline,after row=\hline},
  every last row/.style={after row=\hline}
]{../hw4_code/output/effect_table.csv}%
}
\end{table}

The estimated covariance matrices for each treatment group
are reported in \cref{tab:covariance-matrix-budesonide},
\cref{tab:covariance-matrix-nedocromil}, and
\cref{tab:covariance-matrix-placebo}.

\begin{table}[htbp]
  \centering
  \caption{Estimated covariance matrix under the final model (Budesonide)}
  \label{tab:covariance-matrix-budesonide}
  \begingroup\let\small\scriptsize\makebox[\textwidth][c]{\small
\begin{tabular}{lrrrrrrrrrrrrrrrr}
\hline
Visit & 0 & 2 & 4 & 12 & 16 & 24 & 28 & 36 & 40 & 48 & 52 & 60 & 72 & 84 & 96 & 108 \\
\hline
0 & 0.234 & 0.235 & 0.230 & 0.245 & 0.238 & 0.240 & 0.244 & 0.245 & 0.240 & 0.255 & 0.246 & 0.248 & 0.230 & 0.210 & 0.196 & 0.182 \\
2 & 0.235 & 0.257 & 0.252 & 0.268 & 0.261 & 0.263 & 0.268 & 0.268 & 0.262 & 0.280 & 0.269 & 0.272 & 0.252 & 0.230 & 0.215 & 0.199 \\
4 & 0.230 & 0.252 & 0.263 & 0.281 & 0.273 & 0.275 & 0.280 & 0.280 & 0.275 & 0.293 & 0.281 & 0.285 & 0.263 & 0.241 & 0.224 & 0.208 \\
12 & 0.245 & 0.268 & 0.281 & 0.324 & 0.314 & 0.317 & 0.323 & 0.323 & 0.317 & 0.338 & 0.324 & 0.328 & 0.304 & 0.278 & 0.259 & 0.240 \\
16 & 0.238 & 0.261 & 0.273 & 0.314 & 0.324 & 0.327 & 0.333 & 0.333 & 0.327 & 0.348 & 0.335 & 0.338 & 0.313 & 0.286 & 0.267 & 0.248 \\
24 & 0.240 & 0.263 & 0.275 & 0.317 & 0.327 & 0.355 & 0.362 & 0.362 & 0.355 & 0.378 & 0.364 & 0.368 & 0.341 & 0.311 & 0.290 & 0.269 \\
28 & 0.244 & 0.268 & 0.280 & 0.323 & 0.333 & 0.362 & 0.400 & 0.400 & 0.392 & 0.418 & 0.402 & 0.407 & 0.377 & 0.344 & 0.321 & 0.298 \\
36 & 0.245 & 0.268 & 0.280 & 0.323 & 0.333 & 0.362 & 0.400 & 0.447 & 0.438 & 0.467 & 0.449 & 0.454 & 0.421 & 0.384 & 0.358 & 0.333 \\
40 & 0.240 & 0.262 & 0.275 & 0.317 & 0.327 & 0.355 & 0.392 & 0.438 & 0.468 & 0.499 & 0.480 & 0.485 & 0.449 & 0.411 & 0.383 & 0.355 \\
48 & 0.255 & 0.280 & 0.293 & 0.338 & 0.348 & 0.378 & 0.418 & 0.467 & 0.499 & 0.586 & 0.563 & 0.570 & 0.527 & 0.482 & 0.449 & 0.417 \\
52 & 0.246 & 0.269 & 0.281 & 0.324 & 0.335 & 0.364 & 0.402 & 0.449 & 0.480 & 0.563 & 0.567 & 0.573 & 0.531 & 0.485 & 0.452 & 0.420 \\
60 & 0.248 & 0.272 & 0.285 & 0.328 & 0.338 & 0.368 & 0.407 & 0.454 & 0.485 & 0.570 & 0.573 & 0.626 & 0.580 & 0.529 & 0.494 & 0.458 \\
72 & 0.230 & 0.252 & 0.263 & 0.304 & 0.313 & 0.341 & 0.377 & 0.421 & 0.449 & 0.527 & 0.531 & 0.580 & 0.597 & 0.545 & 0.509 & 0.472 \\
84 & 0.210 & 0.230 & 0.241 & 0.278 & 0.286 & 0.311 & 0.344 & 0.384 & 0.411 & 0.482 & 0.485 & 0.529 & 0.545 & 0.580 & 0.541 & 0.502 \\
96 & 0.196 & 0.215 & 0.224 & 0.259 & 0.267 & 0.290 & 0.321 & 0.358 & 0.383 & 0.449 & 0.452 & 0.494 & 0.509 & 0.541 & 0.574 & 0.533 \\
108 & 0.182 & 0.199 & 0.208 & 0.240 & 0.248 & 0.269 & 0.298 & 0.333 & 0.355 & 0.417 & 0.420 & 0.458 & 0.472 & 0.502 & 0.533 & 0.575 \\
\hline
\end{tabular}
}\endgroup
\end{table}

\begin{table}[htbp]
  \centering
  \caption{Estimated covariance matrix under the final model (Nedocromil)}
  \label{tab:covariance-matrix-nedocromil}
  \begingroup\let\small\scriptsize\makebox[\textwidth][c]{\small
\begin{tabular}{lrrrrrrrrrrrrrrrr}
\hline
Visit & 0 & 2 & 4 & 12 & 16 & 24 & 28 & 36 & 40 & 48 & 52 & 60 & 72 & 84 & 96 & 108 \\
\hline
0 & 0.268 & 0.250 & 0.244 & 0.262 & 0.258 & 0.285 & 0.289 & 0.300 & 0.302 & 0.311 & 0.314 & 0.305 & 0.277 & 0.255 & 0.222 & 0.217 \\
2 & 0.250 & 0.254 & 0.248 & 0.266 & 0.262 & 0.290 & 0.293 & 0.305 & 0.307 & 0.317 & 0.320 & 0.311 & 0.282 & 0.260 & 0.226 & 0.221 \\
4 & 0.244 & 0.248 & 0.262 & 0.280 & 0.276 & 0.306 & 0.309 & 0.321 & 0.324 & 0.334 & 0.337 & 0.327 & 0.297 & 0.274 & 0.238 & 0.233 \\
12 & 0.262 & 0.266 & 0.280 & 0.323 & 0.319 & 0.353 & 0.357 & 0.370 & 0.374 & 0.385 & 0.389 & 0.377 & 0.343 & 0.316 & 0.275 & 0.268 \\
16 & 0.258 & 0.262 & 0.276 & 0.319 & 0.333 & 0.369 & 0.373 & 0.387 & 0.391 & 0.402 & 0.406 & 0.395 & 0.358 & 0.330 & 0.287 & 0.280 \\
24 & 0.285 & 0.290 & 0.306 & 0.353 & 0.369 & 0.444 & 0.449 & 0.466 & 0.470 & 0.484 & 0.489 & 0.475 & 0.431 & 0.397 & 0.346 & 0.337 \\
28 & 0.289 & 0.293 & 0.309 & 0.357 & 0.373 & 0.449 & 0.472 & 0.490 & 0.494 & 0.509 & 0.514 & 0.499 & 0.454 & 0.418 & 0.363 & 0.355 \\
36 & 0.300 & 0.305 & 0.321 & 0.370 & 0.387 & 0.466 & 0.490 & 0.546 & 0.551 & 0.567 & 0.573 & 0.556 & 0.506 & 0.465 & 0.405 & 0.395 \\
40 & 0.302 & 0.307 & 0.324 & 0.374 & 0.391 & 0.470 & 0.494 & 0.551 & 0.582 & 0.599 & 0.605 & 0.588 & 0.534 & 0.491 & 0.428 & 0.417 \\
48 & 0.311 & 0.317 & 0.334 & 0.385 & 0.402 & 0.484 & 0.509 & 0.567 & 0.599 & 0.662 & 0.668 & 0.649 & 0.590 & 0.543 & 0.472 & 0.461 \\
52 & 0.314 & 0.320 & 0.337 & 0.389 & 0.406 & 0.489 & 0.514 & 0.573 & 0.605 & 0.668 & 0.722 & 0.702 & 0.637 & 0.587 & 0.511 & 0.498 \\
60 & 0.305 & 0.311 & 0.327 & 0.377 & 0.395 & 0.475 & 0.499 & 0.556 & 0.588 & 0.649 & 0.702 & 0.727 & 0.660 & 0.608 & 0.529 & 0.516 \\
72 & 0.277 & 0.282 & 0.297 & 0.343 & 0.358 & 0.431 & 0.454 & 0.506 & 0.534 & 0.590 & 0.637 & 0.660 & 0.677 & 0.623 & 0.542 & 0.529 \\
84 & 0.255 & 0.260 & 0.274 & 0.316 & 0.330 & 0.397 & 0.418 & 0.465 & 0.491 & 0.543 & 0.587 & 0.608 & 0.623 & 0.648 & 0.564 & 0.551 \\
96 & 0.222 & 0.226 & 0.238 & 0.275 & 0.287 & 0.346 & 0.363 & 0.405 & 0.428 & 0.472 & 0.511 & 0.529 & 0.542 & 0.564 & 0.567 & 0.553 \\
108 & 0.217 & 0.221 & 0.233 & 0.268 & 0.280 & 0.337 & 0.355 & 0.395 & 0.417 & 0.461 & 0.498 & 0.516 & 0.529 & 0.551 & 0.553 & 0.613 \\
\hline
\end{tabular}
}\endgroup
\end{table}

\begin{table}[htbp]
  \centering
  \caption{Estimated covariance matrix under the final model (Placebo)}
  \label{tab:covariance-matrix-placebo}
  \begingroup\let\small\scriptsize\makebox[\textwidth][c]{\small
\begin{tabular}{lrrrrrrrrrrrrrrrr}
\hline
Visit & 0 & 2 & 4 & 12 & 16 & 24 & 28 & 36 & 40 & 48 & 52 & 60 & 72 & 84 & 96 & 108 \\
\hline
0 & 0.279 & 0.271 & 0.268 & 0.288 & 0.295 & 0.300 & 0.299 & 0.315 & 0.315 & 0.320 & 0.316 & 0.307 & 0.294 & 0.265 & 0.239 & 0.209 \\
2 & 0.271 & 0.278 & 0.275 & 0.296 & 0.302 & 0.309 & 0.307 & 0.323 & 0.323 & 0.328 & 0.325 & 0.315 & 0.302 & 0.273 & 0.246 & 0.215 \\
4 & 0.268 & 0.275 & 0.290 & 0.312 & 0.319 & 0.325 & 0.324 & 0.341 & 0.341 & 0.346 & 0.342 & 0.332 & 0.318 & 0.287 & 0.259 & 0.227 \\
12 & 0.288 & 0.296 & 0.312 & 0.367 & 0.374 & 0.382 & 0.380 & 0.400 & 0.400 & 0.406 & 0.402 & 0.390 & 0.374 & 0.337 & 0.304 & 0.266 \\
16 & 0.295 & 0.302 & 0.319 & 0.374 & 0.413 & 0.421 & 0.420 & 0.442 & 0.442 & 0.448 & 0.443 & 0.430 & 0.412 & 0.372 & 0.335 & 0.294 \\
24 & 0.300 & 0.309 & 0.325 & 0.382 & 0.421 & 0.467 & 0.465 & 0.490 & 0.490 & 0.497 & 0.492 & 0.477 & 0.457 & 0.413 & 0.372 & 0.326 \\
28 & 0.299 & 0.307 & 0.324 & 0.380 & 0.420 & 0.465 & 0.499 & 0.525 & 0.525 & 0.533 & 0.527 & 0.512 & 0.490 & 0.442 & 0.399 & 0.349 \\
36 & 0.315 & 0.323 & 0.341 & 0.400 & 0.442 & 0.490 & 0.525 & 0.592 & 0.592 & 0.601 & 0.595 & 0.577 & 0.553 & 0.499 & 0.450 & 0.394 \\
40 & 0.315 & 0.323 & 0.341 & 0.400 & 0.442 & 0.490 & 0.525 & 0.592 & 0.625 & 0.634 & 0.627 & 0.609 & 0.583 & 0.526 & 0.474 & 0.415 \\
48 & 0.320 & 0.328 & 0.346 & 0.406 & 0.448 & 0.497 & 0.533 & 0.601 & 0.634 & 0.686 & 0.678 & 0.658 & 0.631 & 0.569 & 0.513 & 0.449 \\
52 & 0.316 & 0.325 & 0.342 & 0.402 & 0.443 & 0.492 & 0.527 & 0.595 & 0.627 & 0.678 & 0.717 & 0.696 & 0.667 & 0.602 & 0.542 & 0.475 \\
60 & 0.307 & 0.315 & 0.332 & 0.390 & 0.430 & 0.477 & 0.512 & 0.577 & 0.609 & 0.658 & 0.696 & 0.717 & 0.686 & 0.620 & 0.558 & 0.489 \\
72 & 0.294 & 0.302 & 0.318 & 0.374 & 0.412 & 0.457 & 0.490 & 0.553 & 0.583 & 0.631 & 0.667 & 0.686 & 0.710 & 0.641 & 0.578 & 0.506 \\
84 & 0.265 & 0.273 & 0.287 & 0.337 & 0.372 & 0.413 & 0.442 & 0.499 & 0.526 & 0.569 & 0.602 & 0.620 & 0.641 & 0.652 & 0.587 & 0.514 \\
96 & 0.239 & 0.246 & 0.259 & 0.304 & 0.335 & 0.372 & 0.399 & 0.450 & 0.474 & 0.513 & 0.542 & 0.558 & 0.578 & 0.587 & 0.622 & 0.544 \\
108 & 0.209 & 0.215 & 0.227 & 0.266 & 0.294 & 0.326 & 0.349 & 0.394 & 0.415 & 0.449 & 0.475 & 0.489 & 0.506 & 0.514 & 0.544 & 0.569 \\
\hline
\end{tabular}
}\endgroup
\end{table}

\subsubsection{Hypothesis testing}

Instead of interpreting the key parameters, we performed a test of
parallelism, with 
$H_0: \beta_{visit \times bud} = \beta_{visit \times ned} = \beta_{visit \times plbo}$
and $\beta_{visit \times bud}^{sp} = \beta_{visit \times ned}^{sp} = \beta_{visit \times plbo}^{sp}$
versus $H_1: \text{at least one of these equalities fails}$.
The F statistic is 2.87, and the p-value is 0.022,
which indicates a significant result. Thus we reject
the null hypothesis and conclude that mean PostBD FEV1 changes
differently over visit time among the three treatment groups.

\section{Visualization and Interpretation}

\subsection{Visualization}

The observed PostBD FEV1 means along with fitted PostBD FEV1 means
for each treatment group over time are illustrated
in \cref{fig:fitted-observed-mean-compare}. 
The handling of covariates was consistent with
that in Assigment 3.
The model shows an adequate fit to the data.

\begin{figure}[htbp]
  \centering
  \includegraphics[width=0.9\textwidth]{../hw4_code/output/fitted_observed_mean_compare.png}
  \caption{Observed and fitted mean trajectories by treatment group.}
  \label{fig:fitted-observed-mean-compare}
\end{figure}



\subsection{Summary}

First, we compared five candidate correlation structures
using a saturated mean model with covariate adjustment
and homogneneous variance. We selected ANTE based
on AIC. Then we compared four variance-covariance structures,
fixing ANTE and the mean structure. We selected
heterogeneous variance over time with
group-specific ANTE covariance matrix based on AIC.
Next, fixing the covariance structure,
we compared four
candidate mean structures. We selected the
linear spline model (knot = 60) based on AIC, BIC,
and interpretability. Then we verified that 
the covariance structure we chose initially
continued to provide the most adequate fit
under the linear spline model.

The final model is a linear spline model
with group-specific heterogeneous ANTE covariance matrix.
We performed an F-test of parallelism,
and we find that mean PostBD FEV1
changes differently over time among
the three treatment groups (p = 0.022).

\section{Code Availability}

The code is available at
\url{https://github.com/CHENGHAO-WANG/767_homework/tree/main/hw4_code}.

\end{document}
